\Needspace{9\baselineskip}\begin{center}\itshape Further investigations on the periods of the roots.\\ Sums with an even number of terms are real quantities.\end{center}


\Needspace{6\baselineskip}\Article{355}
Since \(n\) is always assumed to be odd, \(2\) is among the factors of \(n-1\), and the collection \(\Omega\) is formed from \(\frac12(n-1)\) periods of two terms. A period such as \((2,\lambda)\) will consist of the roots \([\lambda]\) and \([\lambda g^{\frac12(n-1)}]\), where, as above, \(g\) denotes any primitive root modulo \(n\). But we have \(g^{\frac12(n-1)}\equiv-1\ (\mathrm{mod.}\,n)\), and therefore \(\lambda g^{\frac12(n-1)}\equiv-\lambda\) (see Article 62), whence \([\lambda g^{\frac12(n-1)}]=[-\lambda]\). Thus, if we suppose \([\lambda]=\cos\frac{kP}{n}+i\sin\frac{kP}{n}\), and therefore \([-\lambda]=\cos\frac{kP}{n}-i\sin\frac{kP}{n}\), the sum is \((2,\lambda)=2\cos\frac{kP}{n}\). From this we draw here only the conclusion that the value of any sum of two terms is a real quantity. Since any period whose number of terms is even \(=2a\) can be divided into \(a\) periods of two terms each, it is clear, more generally, that the value of any sum
\SourceLeaf{446}
with an even number of terms is always a real quantity. Thus, if in Article 352 the factor two is reserved for the last place among the factors \(\alpha,\allowbreak \beta,\allowbreak \gamma\), etc., all operations, until sums of two terms are reached, will be carried out with real quantities. Imaginary quantities will be introduced only when you pass from these sums to the roots themselves.

\Needspace{9\baselineskip}\begin{center}\itshape On the equation that determines the distribution of the roots \(\Omega\) into two periods.\end{center}


\Needspace{6\baselineskip}\Article{356}
The auxiliary equations that determine, for each value of \(n\), the sums constituting the collection \(\Omega\) deserve the closest attention: they are connected in a remarkable way with the deepest properties of the number \(n\). Here, however, we will restrict our investigation to the following two cases. \textit{First}, we will discuss the quadratic equation whose roots are the sums of \(\frac12(n-1)\) terms; \textit{second}, for the case where \(n-1\) contains the factor \(3\), the cubic equation whose roots are the sums of \(\frac13(n-1)\) terms.

Writing, for brevity, \(m\) for \(\frac12(n-1)\), and denoting by \(g\) any primitive root modulo \(n\), the collection \(\Omega\) will consist of two periods \((m,1)\) and \((m,g)\). The first will contain the roots \([1],\allowbreak [gg],\allowbreak [g^4]\ldots[g^{n-3}]\), the second \([g],\allowbreak [g^3],\allowbreak [g^5]\ldots[g^{n-2}]\). Suppose the least positive residues of the numbers \(gg,\allowbreak g^4\ldots g^{n-3}\) modulo \(n\) are, in any order, \(R,\allowbreak R',\allowbreak R''\), etc., and likewise those of \(g,\allowbreak g^3,\allowbreak g^5\ldots g^{n-2}\) are \(N,\allowbreak N',\allowbreak N''\), etc. Then the roots constituting \((m,1)\) will coincide with \([1],\allowbreak [R],\allowbreak [R'],\allowbreak [R'']\), etc., and the roots of the period \((m,g)\) with \([N],\allowbreak [N'],\allowbreak [N'']\), etc. It is now clear that all the numbers \(1,\allowbreak R,\allowbreak R',\allowbreak R''\), etc., are \textit{quadratic residues} of the number \(n\). Since they are all distinct and smaller than \(n\), and their number is \(=\frac12(n-1)\), hence equal to the number of all the positive residues of \(n\) below \(n\), these residues will coincide exactly with those numbers. It follows at once that all the numbers \(N,\allowbreak N',\allowbreak N''\), etc., which are distinct both from one another and from \(1,\allowbreak R,\allowbreak R'\), etc., and together with these exhaust all the numbers \(1,\allowbreak 2,\allowbreak 3\ldots n-1\), must coincide with all the positive \textit{quadratic nonresidues} of \(n\) below \(n\). If we now suppose that the equation whose roots are the sums \((m,1),\allowbreak (m,g)\) is \[xx-Ax+B=0\] then \[A=(m,1)+(m,g)=-1,\quad B=(m,1)\times(m,g)\]
\SourceLeaf{447}

By Article 345, the product of \((m,1)\) and \((m,g)\) is \[=(m,N+1)+(m,N'+1)+(m,N''+1)+\text{etc.}=W\] and can therefore be reduced to a form such as \(\alpha(m,0)+\beta(m,1)+\gamma(m,g)\). To determine the coefficients \(\alpha,\allowbreak \beta,\allowbreak \gamma\), we observe, \textit{first}, that \(\alpha+\beta+\gamma=m\) (since the number of sums in \(W\) is \(=m\)); \textit{second}, that \(\beta=\gamma\) (this follows from Article 350, since the product \((m,1)\times(m,g)\) is a symmetric function of the sums \((m,1),\allowbreak (m,g)\), of which the larger sum \((n-1,1)\) is composed); \textit{third}, since all the numbers \(N+1,\allowbreak N'+1,\allowbreak N''+1\), etc., lie strictly between the limits \(2\) and \(n+1\), it is clear that either no sum in \(W\) reduces to \((m,0)\), and therefore \(\alpha=0\), when among the numbers \(N,\allowbreak N',\allowbreak N''\), etc., \(n-1\) does not occur; or one does, namely \((m,n)\), and hence \(\alpha=1\), when \(n-1\) is found among the numbers \(N,\allowbreak N',\allowbreak N''\), etc. We conclude that in the first case \(\alpha=0,\allowbreak \ \beta=\gamma=\frac12m\), in the second \(\alpha=1,\allowbreak \ \beta=\gamma=\frac12(m-1)\). At the same time it follows, since the numbers \(\beta\) and \(\gamma\) must necessarily be integers, that the first case occurs, that is, \(n-1\) (or, equivalently, \(-1\)) is not found among the nonresidues of \(n\), when \(m\) is even, or equivalently \(n\) is of the form \(4k+1\); the second case occurs, that is, \(n-1\) or \(-1\) is found among the nonresidues of \(n\), whenever \(m\) is odd, or equivalently \(n\) is of the form \(4k+3\)\HistoricalNote{*)}{In this way we have obtained a new proof of the theorem that \(-1\) is a residue of all prime numbers of the form \(4k+1\), and a nonresidue of all those of the form \(4k+3\), which was already proved above (Articles 108, 109 and 262) in several different ways. If we prefer to assume this theorem, there is no need, in distinguishing the two different cases, to take account of the condition that \(\beta,\allowbreak \gamma\) are already integers in themselves.}. Thus the required product, since \((m,0)=m\), \((m,1)+(m,g)=-1\), is in the first case \(=-\frac12m\), in the second \(=\frac12(m+1)\). The required equation is therefore, in the former case, \(xx+x-\frac14(n-1)=0\), whose roots are \(-\frac12\pm\frac12\surd n\); in the latter, \(xx+x+\frac14(n+1)=0\), whose roots are \(-\frac12\pm\frac12i\surd n\).

Thus, whichever root from \(\Omega\) is taken for \([1]\), the difference between the sums \(\Sigma[\mathfrak R]\) and \(\Sigma[\mathfrak N]\), where for \(\mathfrak R\) we substitute all the positive quadratic residues and for \(\mathfrak N\) all the positive quadratic nonresidues of \(n\) below \(n\), will be \(=\pm\surd n\) for \(n\equiv1\), and \(=\pm i\surd n\) for \(n\equiv3\ (\mathrm{mod.}\,4)\). It also follows easily that, if \(k\) denotes any integer not divisible by \(n\), we have \[\begin{gathered}\Sigma\cos\frac{k\mathfrak RP}{n}-\Sigma\cos\frac{k\mathfrak NP}{n}=\pm\surd n\quad\text{and}\\\Sigma\sin\frac{k\mathfrak RP}{n}-\Sigma\sin\frac{k\mathfrak NP}{n}=0\end{gathered}\] for \(n\equiv1\ (\mathrm{mod.}\,4)\); on the other hand, for \(n\equiv3\ (\mathrm{mod.}\,4)\) the former difference is \(=0\), the latter
\SourceLeaf{448}
\(=\pm\surd n\). These theorems are especially worth remembering for their elegance. We also observe that the upper signs always hold when for \(k\) we take unity or, more generally, a quadratic residue of \(n\), and the lower signs when for \(k\) we take a nonresidue; also that these theorems can be extended to any composite values of \(n\) with their elegance preserved, or rather increased. But we must refrain from discussing these matters here, since they require a deeper investigation, and reserve their consideration for another occasion

